\begin{lemma}[Partial differentiation drops total degree]
Let $\sigma$ be any index type, let $\varphi\in\mathbb{R}[X_a:a\in\sigma]$ and
let $i\in\sigma$. Then
\[
  \deg(\partial_i\varphi)\ \le\ \deg\varphi-1,
\]
where $\deg$ denotes total degree and the subtraction is truncated at $0$
(i.e. $\deg\varphi-1$ means $0$ when $\deg\varphi=0$).
\end{lemma}
\begin{proof}
If $\partial_i\varphi=0$ then $\deg(\partial_i\varphi)=0$ and the inequality is
immediate, since the right-hand side is a natural number. So assume
$\partial_i\varphi\ne0$. The total degree of a nonzero polynomial is the maximum
of $\sum_a s_a$ over the exponent tuples $s$ in its support, so it suffices to
show
\begin{equation}\label{eq:tdp-goal}
  \sum_a s_a\ \le\ \deg\varphi-1
  \qquad\text{for every }s\in\operatorname{supp}(\partial_i\varphi).
\end{equation}

Fix such an $s$, so $\operatorname{coeff}_s(\partial_i\varphi)\ne0$. Let
$\delta_i\in\mathbb{N}^{\sigma}$ be the tuple with $i$-th entry $1$ and all
other entries $0$. The coefficients of a formal partial derivative are given by
\[
  \operatorname{coeff}_s\bigl(\partial_i\varphi\bigr)
  =\operatorname{coeff}_{s+\delta_i}(\varphi)\cdot\bigl(s_i+1\bigr),
\]
which is the defining formula for $\partial_i$ read coefficientwise. Since the
left-hand side is nonzero, the factor $\operatorname{coeff}_{s+\delta_i}(\varphi)$
must be nonzero as well, i.e.
\[
  s+\delta_i\in\operatorname{supp}\varphi .
\]
Consequently, by the background proposition bounding the exponent sums on the
support by the total degree,
\begin{equation}\label{eq:tdp-bound}
  \sum_a\bigl(s+\delta_i\bigr)_a\ \le\ \deg\varphi .
\end{equation}
Finally the exponent sums are related by
\[
  \sum_a\bigl(s+\delta_i\bigr)_a=\sum_a s_a+\sum_a(\delta_i)_a=\sum_a s_a+1,
\]
because the two tuples are added entrywise and $\delta_i$ has entry sum $1$.
Substituting this into \eqref{eq:tdp-bound} gives
$\sum_a s_a+1\le\deg\varphi$, hence $\sum_a s_a\le\deg\varphi-1$ (and the
truncated subtraction causes no difficulty, since $\deg\varphi\ge1$ here). This
is \eqref{eq:tdp-goal}, and the lemma follows.
\end{proof}
